\chapter{Installation}
\label{kap:installation}

ESIM wird aus dem Quelltext gebaut; für Windows gibt es zusätzlich ein
fertiges Paket aus der automatischen Übersetzung auf GitHub. Der Simulator
\esim\ braucht nur einen C++17-Compiler und CMake, der Editor zusätzlich
Qt~6.

\section{Was gebraucht wird}

\begin{center}
\begin{tabularx}{\textwidth}{@{}lLl@{}}
  \toprule
  Werkzeug & Zweck & nötig für \\
  \midrule
  CMake $\geq$ 3.18 & Build-System & alles \\
  C++17-Compiler & Clang (macOS), MSVC (Windows), GCC (Linux) & alles \\
  Eigen 3, nlohmann/json & Bibliotheken; fehlen sie, lädt CMake sie selbst herunter & alles \\
  Qt $\geq$ 6.5 & Oberfläche des Editors & \editor \\
  Python 3 & Testsuite, Werkzeuge, Python-Schnittstelle & optional \\
  numpy, matplotlib & Ergebnisse als Arrays, Plotten mit \datei{plot.py} & optional \\
  ngspice & drei Vergleichstests gegen ngspice & optional \\
  \bottomrule
\end{tabularx}
\end{center}

Ohne Qt baut CMake nur den Simulator und seine Tests; alles im Teil III
dieses Handbuchs funktioniert dann trotzdem über die Kommandozeile.

\section{macOS}

Mit Homebrew sind alle Voraussetzungen schnell installiert:

\begin{shell}
brew install cmake qt eigen nlohmann-json
\end{shell}

Dann im Ordner des Quelltexts:

\begin{shell}
cmake -S . -B build
cmake --build build --parallel
ctest --test-dir build                 # Tests (optional)
\end{shell}

Der Editor liegt danach als Programmpaket in
\datei{build/editor/esim-editor.app}:

\begin{shell}
open build/editor/esim-editor.app
\end{shell}

Ein eigenständiges Programmpaket, das man in den Ordner »Programme« ziehen
kann (mit den Qt-Bibliotheken, \esim\ und den Beispielen im Paket), legt

\begin{shell}
cmake --install build --prefix dist
\end{shell}

in \datei{dist/esim-editor.app} an.

\begin{hinweis}[Gatekeeper]
  Das Programm ist nicht mit einer Apple-Entwickler-ID signiert. Ein selbst
  gebautes Programm startet ohne Rückfrage. Wird das Paket auf einen
  anderen Mac kopiert, verweigert macOS beim ersten Start unter Umständen das
  Öffnen; dann im Finder mit Rechtsklick \pfad\ \menue{Öffnen} starten oder
  unter \menue{Systemeinstellungen}\pfad\menue{Datenschutz \& Sicherheit}
  einmal erlauben.
\end{hinweis}

\section{Windows 11}

\subsection{Fertiges Paket}

Bei jedem Hochladen auf GitHub baut GitHub Actions ein Windows-Paket. Es
liegt im Repository unter \menue{Actions}, im jeweils neuesten Lauf unter
\menue{Artifacts}, als \datei{esim-windows}. Alternativ mit dem
GitHub-Kommandozeilenprogramm:

\begin{shell}
gh run list --limit 3                    # Nummer des neuesten Laufs
gh run download <Nummer> -n esim-windows
\end{shell}

Das Paket entpacken und \datei{esim-editor.exe} starten. Qt muss dafür nicht
installiert sein; die Qt-Bibliotheken, \datei{esim.exe} und die Beispiele
liegen im selben Ordner. Ein Installationsprogramm gibt es nicht – der
Ordner kann an eine beliebige Stelle kopiert werden.

\begin{hinweis}[SmartScreen]
  Auch unter Windows ist das Programm nicht signiert. Beim ersten Start kann
  SmartScreen warnen; mit \menue{Weitere Informationen}\pfad\menue{Trotzdem
  ausführen} lässt es sich starten.
\end{hinweis}

\subsection{Selbst bauen}

Benötigt werden Visual Studio 2022 oder neuer (oder nur dessen Build Tools)
mit der Arbeitslast »Desktopentwicklung mit C++«, CMake, Git, Qt~6 für MSVC
und optional Python. Mit \code{winget} geht das so:

\begin{shell}
winget install --id Git.Git -e
winget install --id Kitware.CMake -e
winget install --id Python.Python.3.12 -e
winget install --id Microsoft.VisualStudio.2022.BuildTools -e --override "--wait --passive --add Microsoft.VisualStudio.Workload.VCTools --includeRecommended"
python -m pip install aqtinstall
python -m aqt install-qt windows desktop 6.8.3 win64_msvc2022_64 -O C:\Qt
\end{shell}

Dann in einer neuen Eingabeaufforderung im Ordner des Quelltexts:

\begin{shell}
cmake -S . -B build -DCMAKE_PREFIX_PATH=C:/Qt/6.8.3/msvc2022_64
cmake --build build --config Release --parallel
ctest --test-dir build -C Release --output-on-failure
cmake --install build --config Release --prefix dist
\end{shell}

\texttt{dist\textbackslash esim-editor.exe} läuft danach für sich allein.
CMake wählt das neueste installierte Visual Studio von selbst; Qts
Bibliotheken für MSVC 2022 passen auch zu neueren Versionen.

\begin{achtung}[Visual Studio 2026]
  Visual Studio 2026 wird erst von CMake 4.2 an erkannt. Mit einem älteren
  CMake bricht die Konfiguration ab; dann ein neueres CMake installieren
  oder das mit Visual Studio gelieferte verwenden
  (\texttt{Common7\textbackslash\allowbreak IDE\textbackslash CommonExtensions\textbackslash Microsoft\textbackslash CMake\textbackslash CMake\textbackslash bin}).
\end{achtung}

Alle Programme arbeiten intern mit UTF-8; Dateien und Ordner mit Umlauten
(etwa \texttt{C:\textbackslash\allowbreak Übung\textbackslash\allowbreak Schaltkreis~ä.json}) sind kein Problem.

\section{Linux}

Der Simulator und der Editor bauen auch mit GCC unter Linux; ein Paket wird
dafür nicht gepflegt. Mit den Paketen der Distribution (z.\,B. unter Debian
und Ubuntu \code{cmake g++ qt6-base-dev libeigen3-dev nlohmann-json3-dev})
genügen dieselben Befehle wie unter macOS.

\section{Optionale Teile}

\paragraph{Python.}
Die Python-Schnittstelle (Kapitel~\ref{kap:python}) braucht nur die
Standardbibliothek. Mit \code{pip install numpy matplotlib} liefert sie
numpy-Arrays, und \datei{tools/plot.py} kann zeichnen.

\paragraph{ngspice.}
Ist ngspice installiert (unter Windows \datei{ngspice_con.exe} im Suchpfad),
legt CMake drei zusätzliche Tests an, die esims Modelle (BSIM3, MOS Level~3,
SPICE-Teilschaltungen) gegen ngspice vergleichen. Für die Benutzung ist
ngspice nicht nötig.

\section{Prüfen, ob alles läuft}

\begin{shell}
./build/esim examples/rc_lowpass.json -o results
\end{shell}

Auf der Konsole erscheinen der Arbeitspunkt und eine Zeile je Analyse; die
Kurven liegen danach in \datei{results/rc_lowpass.ac.csv} und
\datei{results/rc_lowpass.tran.csv}. Im Editor öffnet
\menue{Datei}\pfad\menue{Beispiele} dieselbe Datei; \tasten{Strg}{R} rechnet
sie.
